\subsection{The group representation and its fixed menus}

Let $X=(u,\beta)$, with $u$ of weight one and $\beta$ of weight $j$.
As actual site data in the buffered chart, $\beta$ is the uniquely
determined lift integer.  As an argument of a polynomial identity,
it is a formal variable.

\paragraph{Exact marked projections.}\label{ap:marking-conventions}
A rational homomorphism $\pi:G\to Y$ between filtered groups is
\emph{strong} if $\pi(G_i)=Y_i$ at every layer. A \emph{marked}
polynomial map is supplied with a fixed polynomial $P_Y$ and an
exact identity $\pi(g(X))=P_Y(X)$ in its formal variables $X$.
Retaining that identity preserves the old determining equations
when a new test is constructed. The rational refiltration statement
that preserves this structure is Lemma~\ref{lifting:refiltration}.

Rational kernels, simultaneous intersections, sections, controlled
lattices and fixed-step Baker--Campbell--Hausdorff calculations have
polynomial logarithmic cost by Lemma~\ref{prelim:rational-operations}.
All uses of this calculus allow replacement of the degree by any
fixed larger degree; none bounds arbitrary orbit coefficients.

\begin{lemma}[Representation with exact determining projection]
\label{ap:ascent:representation}
At a fixed parent and feasible center, the ascent data admit the
following representation, after pigeonholing warm structural data.

At a passive layer $j>s$, there is one triangular shear group $G^0$
for the absolute slots, of degree at most $s$, and path-dependent
ordinary filtered polynomial maps $g_\psi(t)$ into it.

At an active layer $j\leq s$, there is one joint triangular group
$G^0$ for the active determining slots $j<i\leq s$ and the absolute
slots, together with a strong rational projection
$\pi^0:G^0\to Y$ to the determining shear group.
There is a fixed
weighted filtered polynomial $P_Y(X)$ such that
\begin{equation}\label{ap:ascent:marking}
 \pi^0g_\psi(t)=P_Y(X_\psi(t)),\qquad
 X_\psi(t)=(\psi(t),\beta(\psi(t))).
\end{equation}

In both cases the two payoffs, excluding the slice, current flag
and target multiplier, are evaluations of fields
\[
 F^-_{\psi,u},F^+_{\psi,u}:G^0/\Gamma^0\longrightarrow[0,1].
\]
They have a common warm logarithmic Lipschitz bound.  At every
specified accuracy they have meshes of polynomial logarithmic size
in two norms.  A point-value mesh approximates the nilmanifold
functions $F^\pm_{\psi,u}$ in supremum norm on $G^0/\Gamma^0$.
A field mesh approximates $u\mapsto F^\pm_{\psi,u}$ simultaneously
over all parent sites, with the supremum also over $G^0/\Gamma^0$.
None of these bounds depends on orbit coefficient sizes.  The original
cutoff functions may be retained exactly in the output; the meshes
are only devices for selecting partners and comparing scores.
\end{lemma}

\begin{proof}
We first separate the passive determining conditions from the group
variables.  Outside the required current chart choose fixed common
values, which the local and forecasting supports never test.
At each site in the chart use fixed padded boxes, for example radii
$2w_i$ around the prepared offsets, to compute the later old slots
successively.  Triangularity gives at most one tuple.  If a required
padded lift does not exist, set the resulting passive multiplier
equal to zero.  Evaluate the passive cutoffs and congruences on this
canonical tuple, obtaining site multipliers $L_\psi^\pm$.

This procedure can use active determining slots, but it does not
alter a nonzero payoff.  On the actual orbit, every nonzero active
cutoff forces exactly its canonical padded lift.  The same is true
on any returned orbit having the exact determining projection
\eqref{ap:ascent:marking}.  Thus the separated expression agrees
with the original joint expression at every point at which it will
be tested, including when either expression is zero.  On arbitrary
points of the full group quotient we simply keep $L_\psi^\pm$ as
site multipliers independent of that group point; no identity with
unconstrained slot evaluation is claimed there.

The canonical tuple is fixed as a function of the site.  Only the
bounded centers, cutoff radii, moduli and classes vary with the
path.  Meshing those real parameters and enumerating the discrete
ones gives a field mesh in site supremum norm.  Its cost is warm
at $j$, since only widths $Q_i$ with $i>j$ enter.  The fixed input
$f_0$ can be included without increasing this mesh size.

For active determining classes write
$b_i=M_i k_i+\alpha_i$ and scale their residuals by $1/M_i$.
Pigeonhole the classes and moduli, as well as the absolute rank and
weight assignments, before fixing the group.  This has warm cost.
Order determining slots before absolute slots.  At the first active
step there are no previously processed active determining slots;
at later steps the larger absolute weights permit this order.

Apply the shear construction in the proof of
Lemma~\ref{counting:conversion}.  In the residual triangular map,
normalize all pure full-destination-weight slot coefficients by
triangular integral changes of the summation variables, in increasing
slot order.  These changes are independent of $t$, preserve the
integer tuples bijectively, and are substituted exactly into every
later formula.  Large coefficients created by an earlier change
are reduced when their own slot is reached.  Write the resulting
residual map as
\[
                           S^0\circ g_\psi(t),
\]
where $S^0$ consists of its full-destination-weight slot terms.
Those coefficients are bounded; its determining part is fixed
at the parent, since full-weight slot coefficients cannot depend
on the positive-weight variables $X$.  The remaining transformations
strictly lower slot weight and form the filtered group $G^0$, with
filtration given by weight deficit.

Forgetting absolute slots is a group homomorphism to $Y$: a
determining transformation has no absolute argument.  A determining
shear of any deficit lifts by acting on determining slots and
leaving absolute slots fixed.  The projection is therefore strong
at every filtration layer, with controlled rational sections.
Taking the determining part of $(S^0)^{-1}$ composed with the exact
residual map defines $P_Y(X)$ and proves
\eqref{ap:ascent:marking} as a polynomial identity.  Cutoff offsets
are placed in the tested functions, not in these target polynomials.
All integral normalizations can be undone when raw determining
coordinates are restored.

The two functions on the quotient evaluate the absolute and active
determining cutoffs, multiplied respectively by
$f_0(u)L^-_\psi(u)$ and $L^+_\psi(u)$.  They exclude the current
flag, slice and $\lambda$.  The integer-lattice action permutes
tuples, so these are well-defined quotient functions.  Triangular
uniqueness gives their cap one.  The proof of
Lemma~\ref{counting:conversion} bounds their Lipschitz constants:
on bounded group representatives, the inverse of the bounded
full-weight map $S^0$ and its derivatives have dimension-controlled
logarithmic bounds; contributing integer tuples and their derivative
estimates have the same bounds.  Class scalings and inherited
cutoff radii have warm cost here.

Finally, $S^0$ and all cutoff parameters lie in bounded,
polynomial-dimensional parameter sets with controlled moduli of
continuity.  Their grids, together with the passive field mesh
already constructed, give both claimed meshes.  No mesh is taken
over the polynomial coefficients of $g_\psi$ or over arbitrary
site functions.  Keeping the original functions after a mesh-based
selection costs only the prescribed approximation error, so their
exact widths and buffers remain available.
\end{proof}

\subsection{Reducing degree while retaining a marked projection}

\subsubsection{Symbols and real or rational control}\label{ap:symbol-conventions}
For a filtered group $G$ of degree $q$, write
$\mathfrak g_i=\log G_i$.  Its associated graded Lie algebra is
$\bigoplus_{i=1}^q\mathfrak g_i/\mathfrak g_{i+1}$.
The \emph{symbol} of $g$ is the map into its simply connected group
whose logarithmic component in grade $i$ is the degree-$i$
homogeneous part of $\log g$ projected to that grade; with assigned
variable weights, homogeneous degree means weighted homogeneous
degree. In the scalar
case, the symbol is just the homogeneous polynomial at the assigned
weight.
On a box with side lengths $H=(H_1,\ldots,H_n)$, write
$H^\nu=\prod_iH_i^{\nu_i}$. A symbol is \emph{slow} at budget
$p$ if its logarithmic coefficient of $u^\nu$ has magnitude at most
$\mathcal B(p)H^{-\nu}$ in the declared coordinates. It is
\emph{rational} if its logarithmic coefficients have a common
denominator at most $\mathcal B(p)$, without a numerator bound.
A polynomial map is \emph{fully slow} if the corresponding bounds
hold for every coefficient of its logarithm, including the constant
coefficient, after translating the box to the origin.
Thus slow control concerns normalized real variables, whereas
rational control concerns denominators in the raw integer variables.
The compatible splittings and their lifts to actual maps are proved
in Section~\ref{prelim:section}.

The group representation separates two requirements.  The functions
must retain a positive signed score, and the orbit must retain the
fixed determining projection.  The algebraic reason these requirements
can coexist is that the projection can record the top coordinate
removed in a degree reduction.
In the passive case take $Y$ to be the trivial group and $\pi^0$
the trivial projection.  Write $r$ for the degree of the represented
group $G^0$.

Suppose a strong projection $\pi:G\to F$ is injective on its top
layer $G_r$.  Then an element $g\in G$ is determined by
$gG_r$ and $\pi(g)$: two elements with the same first coordinate
differ by an element of $G_r$, and equality of their projections
makes that difference trivial.  A prescribed polynomial projection
therefore leaves only a degree-$(r-1)$ orbit to construct.  The
original group need not satisfy this injectivity.  The next lemma
allows us to obtain it on the rational subgroup selected by the
symbol argument.

\begin{lemma}[Refiltration under a strong projection]
\label{lifting:refiltration}
Let $\pi:G\to F$ be a rational homomorphism of connected simply
connected filtered nilpotent Lie groups such that
$\pi(G_i)=F_i$ for every $i$.  Let $\mathfrak v$ be a graded rational
subalgebra of $\operatorname{gr}G$, and let
$\mathfrak v_F=\operatorname{gr}(\pi)(\mathfrak v)$.
Define $H_i\subset G_i$ and $H_{F,i}\subset F_i$ by taking the
preimages of $\mathfrak v_i$ and $(\mathfrak v_F)_i$ in each graded
quotient.  Then $H=H_1$ and $H_F=H_{F,1}$ are filtered rational
groups, and $\pi(H_i)=H_{F,i}$ for every $i$.

If the original degree is $r$ and the map
$\mathfrak v_r\to(\mathfrak v_F)_r$ is injective, then
\begin{equation}
 H\longrightarrow
 (H/H_r)\mathop{\times}_{H_F/H_{F,r}}H_F,
 \qquad h\longmapsto(hH_r,\pi(h))
 \label{lifting:fiber-map}
\end{equation}
is a rational isomorphism, onto at each filtered layer.
\end{lemma}

\begin{proof}
Every $H_i$ contains $G_{i+1}$, so the $H_i$ are nested.
The bracket inclusion follows by reducing a bracket modulo
$G_{i+j+1}$ and using
$[\mathfrak v_i,\mathfrak v_j]\subset\mathfrak v_{i+j}$.
Rationality follows from the rationality of the quotient subspaces.
For $y\in H_{F,i}$ choose $x\in G_i$ whose graded image lies in
$\mathfrak v_i$ and projects to that of $y$.
The discrepancy $\pi(x)^{-1}y$ lies in $F_{i+1}$, and hence has a
lift in $G_{i+1}\subset H_i$.  Correcting $x$ proves strongness.

At the top, $G_{r+1}=F_{r+1}=1$, so $H_r\to H_{F,r}$ is both
injective and surjective.  The kernel of
\eqref{lifting:fiber-map} is therefore trivial.
Given a compatible pair $(xH_r,y)$, compatibility says
$\pi(x)^{-1}y\in H_{F,r}$.  There is a unique $z\in H_r$ with
$\pi(z)=\pi(x)^{-1}y$, and $xz$ maps to the pair.
The same correction lies in every required earlier layer, proving
filtered surjectivity.  In logarithmic coordinates the map and its
inverse are rational linear maps between rational Lie algebras,
followed by exponential coordinates.  This proves rationality and
the claimed isomorphism.
\end{proof}

Here is how the local correlations supply the hypothesis of the lemma.
Truncate the top Fourier expansion of each payoff and retain those
modes whose tested averages are significant relative to the score.
The total of the discarded averages is charged to the error.  Each
retained frequency can be paired, by path detection, with a top
frequency of an ambient partner $Q_i(u)$ in a group $G_{Q_i}$.
Step drop puts the joint symbol in a rational graded algebra on which
the sum of those two frequencies vanishes, up to slow and rational
outer factors.  Choose a basis among the payoff frequencies restricted
to $G^0_r\cap\ker\pi^0$, and let $L$ be their connected joint
annihilator there.  After averaging over $L$, form
\[
 G=(G^0/L)\times\prod_iG_{Q_i},
 \qquad F=Y\times\prod_iG_{Q_i},
\]
with projection induced by $\pi^0$ and the identities on the partners.
Lemma~\ref{prelim:symbol-splitting} intersects the pair constraints in
one graded algebra $\mathfrak v$.

A vector in $\mathfrak v_r$ whose projection to $F$ is zero has
zero partner coordinates.  Each pair constraint then says that the
corresponding payoff frequency vanishes on its $G^0$ coordinate.
That coordinate lies in $L$, so it is zero in $G^0/L$.
Thus $\mathfrak v_r\to(\mathfrak v_F)_r$ is injective, exactly as
required in Lemma~\ref{lifting:refiltration}.  Only a basis of
restricted frequencies is needed, so the number of partner groups is
bounded by the dimension of the top kernel.

The next task is to factor the joint local orbit in $G$ as
\[
 g=EhR,\qquad
 \pi E=E_F,\quad \pi h=P_H,\quad \pi R=R_F,
\]
where the projected factors are fixed ambient polynomials and
$h$ is adapted to the refiltration $H$.  Matching symbols alone
does not give these identities; the proofs below construct exact
kernel corrections.  Once they hold,
\eqref{lifting:fiber-map} reconstructs $h$ from $hH_r$ and the
fixed value $P_H$.  On suitable lattice covers, a tested function
of $EhR$ becomes a function of this compatible pair.  Specializing
its second coordinate at $P_H$ leaves a field of tests on the
lower-degree quotient $H/H_r$.

Proposition~\ref{lifting:candidate-globalization} carries out this
construction when all auxiliary slice and score losses have the same
early budget.  We need two refinements.  At passive layers the losses
in the scoring slices must all be warm; the next lemma recovers that
bound in the symbol-matching step.  At active layers the auxiliary
construction may be cold, so Section~\ref{ap:active} retains its
short variables and outer factors as polynomial labels and preserves
the original score by identities of expected payoffs.

\subsection{Warm matching at a passive layer}

At $j>s$, all degrees that occur in ordinary partner selection are
strictly below $j$.  The warm approximation conclusion of
Theorem~\ref{ap:path:detection} therefore applies.  Its ambient
partners are ordinary niltests times twists localized to the chart.
On a path, their smooth arguments $\psi(t)/H$ and $y(\psi(t))/v$
have bounded normalized polynomial coefficients.  Warm subdivision
and bounded-modulus restriction freeze these twists to any warm
accuracy.  In particular, freezing them does not cost $Q_j$.

The remaining issue is the major calculation in marked-symbol
matching.  The following observation is the reason it has warm
output at passive layers despite using a cold auxiliary theorem.

\begin{lemma}[Warm ordinary matching]\label{ap:passive:warm-matching}
Fix $j>s$ and $1\leq d\leq s$.  Consider polynomially many fixed
weighted homogeneous scalar forms of degree $d$ whose leading path
evaluations have slow-plus-rational decompositions with common warm
bounds on a warm-inverse fraction of paths.  The forms themselves
have slow-plus-rational decompositions on $u$, with coefficients
bounded in $u/H$ and with a common denominator of warm cost.
The auxiliary cold-pre requests can be forecast before the late
path scales and rank thresholds.  They cause no actual cold-cost
restriction of a productive scoring slice.
\end{lemma}

\begin{proof}
A weighted form of degree $d<j$ has no $\beta$ argument.  Apply
Proposition~\ref{lifting:candidate-major}, or its quantitative
refinement in Lemma~\ref{ap:major:refinement}, at cold-pre cost.
It yields, on the ordinary $u$ variables,
\[
                           A(u)=S(u/H)+R(u),
\]
with cold-pre coefficient and denominator bounds.  Any restrictions
on lift-value spaces in that conclusion are irrelevant to an
expression with no lift argument.

For this auxiliary argument retain only the parallel long parameter
variables, whose sides are $L$.  Choose $L$ after the just-obtained
cold-pre output bounds.  The displayed decomposition and the
original warm decomposition have equal evaluations there.  Their
difference is simultaneously slow and rational, with cold-pre
bounds.  Separation, Lemma~\ref{prelim:separation}, makes the
positive-degree difference zero once $L$ exceeds those bounds.
Thus the evaluation of $S$ has the original warm normalized
coefficient bound, and that of $R$ has the original warm denominator
bound.  There is no constant term to separate at this positive
homogeneous grade.

The spatial kernel contains an invertible block at the common warm
normalized scale.  Restricting to this block recovers the coefficients
of $S$ on $u/H$ with a warm inverse norm, giving the claimed real
bound.  For the rational bound, multiply $R$ by the warm common
denominator of its evaluation and apply
Lemma~\ref{ap:major:denominator} to the pure-free spatial
outputs.  Its loss depends on dimensions and the inverse retained
path probability, not on the cold a priori denominator ceiling.
Use the unconditional uniform-goodness estimate, choosing its
exceptional probability smaller than the warm-inverse fraction
under consideration.  This does not assert equidistribution after
conditioning on the productive subset.  At least one retained path is good,
and its denominator conclusion is a warm bound for the fixed $R$.
Apply the same argument to the polynomially many coordinates
jointly, with common denominators throughout.

The enlarged upper ceilings and comparison accuracies are only
cold-pre requests.  They are fixed before the late choices of $L$,
perturbation sizes, rank and parent sides.  At each successive
matching grade the modular loss is chosen from the current warm
probability and accuracy, rather than from an output accuracy not
yet known.  There are only boundedly many grades.  This order avoids
any circular requirement.  Finally, the restriction to parallel
variables above proves an identity of the fixed scalar forms; it
is not a restriction of the scoring slice.  Its long thresholds
therefore do not become cold subslice costs in globalization.
\end{proof}

\subsection{Returning an ordinary orbit}

\begin{proposition}[Passive micro replacement]
\label{ap:passive:replacement}
Suppose $j>s$ and the ascent format holds on a warm-inverse fraction
of paths at a fixed parent and feasible center.  Use the group and
the two fields of Lemma~\ref{ap:ascent:representation}.
There is a single ambient ordinary filtered polynomial
\[
                            g^U(u)\in G^0
\]
of degree at most $s$, with no $\beta$ argument, for which a
warm-inverse fraction of the productive paths have warm-cost
subslices on which replacing $g_\psi$ by $g^U\circ\psi$ preserves
a positive signed full-mass score of warm size, after the allowed
small target discounts.  Both inherited cutoff functions and all
their exact determining conditions are retained.
\end{proposition}

\begin{proof}
We adapt the proof of Proposition~\ref{lifting:candidate-globalization}
and record the points needed for the two payoffs and the warm
budget.  This is a bounded-degree induction inside the original
triangular group.  No new absolute slots are introduced by it.

\emph{Top modes and a fixed menu of inputs.}
Truncate the top vertical Fourier expansions of both fields with
total uniform error much smaller than the current score.  On each
productive path call a restricted top frequency significant if
the absolute value of its tested average in either payoff exceeds
the score divided by a sufficiently large multiple of the total number of retained
modes.  The sum of all discarded averages is negligible.  Take
the joint connected annihilator $L$ in the kernel top torus and
choose a basis among all significant restricted frequencies from
the two lists.  Averaging both fields over $L$ preserves their
range $[0,1]$ and changes the signed score only by the discarded
error.  The number of pivots is at most the top dimension.  As in
the cited proof, pigeonholing bases, rather than arbitrary subsets
of frequencies, costs only a warm fraction.

For each pivot use the payoff in which it is significant.
Path-dependent external fields require care here.  First choose
a simultaneous field mesh at the required site-supremum accuracy.
For every member of that fixed menu form point-value bins using a
common point-value mesh.  The masked site inputs belonging to
these bins now form a fixed controlled menu over all paths.
Apply Theorem~\ref{ap:path:detection} to each bounded masked input,
with error small enough to pay for the menu size, pivot threshold
and retained path probability.  A union bound and Markov's
inequality make the models simultaneously valid on a retained
warm-inverse fraction of paths, for their separately chosen slices
and tests.  The local testing degree is at most $s<j$, so every
output bound in this step is warm.

For clarity, if there are $J$ inputs, retained probability $\kappa$
and required local tolerance $\tau$, choose testing-seminorm errors
at most $\kappa\tau/(4J)$.  If $e_h(\psi)$ is the pathwise supremum
error for input $h$, then
\[
 \Pr\{\max_h e_h(\psi)>\tau\}
 \leq\sum_h\frac{\E e_h(\psi)}{\tau}\leq\kappa/4.
\]
This controls path-dependent tests without attempting to pigeonhole
arbitrary site functions or orbit coefficients.

\emph{Splitting and matching.}
A model term supplies an ambient partner for each pivot.  Truncate
its top modes, retain a biased pairing, and freeze its chart twist
by the warm subdivision described above.  Freeze axes shorter
than the warm step-drop threshold.  Step drop gives pure-long
symbol splittings for the local map and its partners.  Refer the
homogeneous identities back through their affine reindexings;
translation does not affect the leading symbols and residue
dilations cost warm factors.  The separate slices for different
pivots need not meet, since these are homogeneous identities.
Separation joins their fast algebras in the product of the quotient
group and the polynomially many selected partner groups.

The fixed marked polynomial in that product consists initially
only of the ambient partner maps, all functions of $u$ alone.
At a matching grade $d\leq s$, its residual modulo the projected
fast algebra has a warm slow-plus-rational path evaluation.
Lemma~\ref{ap:passive:warm-matching} gives the same bounds for its
ordinary ambient scalar coordinates.  Rational sections lift them
to the required left slow and right rational factors.  Updating
these factors changes no lower grade.  There are at most $s$
grades, so the induction uses boundedly many polynomial increases
of warm budgets.  The argument has no lift argument to project
away: one may take the full lift-value space throughout.  All full
factors and marked maps are ordinary polynomials in $u$.

\emph{Exact reconstruction and field meshes.}
Refilter and reconstruct the lower-degree problem as in
Proposition~\ref{lifting:candidate-globalization}.  Its filtered
rational fiber isomorphism preserves the marked projection
exactly.  Slow kernel factors can be frozen on warm subboxes, and
rational kernel factors on warm residue classes.  Common
denominators, not coefficient numerators, control the latter
restrictions.  The source lattice and cover may therefore be
chosen at warm cost independently of the unrestricted orbit
coefficients.

We verify that the function-menu hypothesis persists under these
operations.  After partners, common marked maps and frozen shifts
have been selected, their site-dependent specializations are fixed
functions of the site.  On the rational joint image used in the
reconstruction proof, a test has the form
\[
             f_{\phi,e,r}(i(h))=\phi(ehr\Gamma).
\]
Here $e$ is bounded and $r$ ranges over finitely many bounded
right-lattice residues.  These functions have one warm Lipschitz
bound.  For a fixed residue $r$, their supremum distance is at most
\[
 \|\phi-\phi'\|_\infty+
            \Lip(\phi')d(e,e').
\]
Use the same clipped McShane extension from the same joint image
for all these tests.  Its infimum formula is nonexpansive in the
supremum norm.  Specialization at the compact marked coordinate
has the common Lipschitz bound proved in the cited reconstruction
argument.  Consequently the original parameter meshes still
control whole fields once the site maps have been fixed.  Point
values have meshes obtained also by gridding that compact marked
coordinate.  No coefficient bound on its site polynomial is
needed.  Equivalently, path-variable cutoff parameters may first
be meshed in field supremum norm and pigeonholed, with the
prescribed error restored after the comparison.

The lower-degree induction now applies to both reconstructed
payoffs at once.  Lifting through the rational fiber isomorphism
and restoring the fixed marked projection give one ordinary
ambient map in the quotient of $G^0$.  A filtered logarithmic
section lifts it to $G^0$.  Its degree remains at most $s$.

\emph{Deaveraging the combined score.}
At every top-torus averaging step, deaverage only after the lower
degree return, and use the \emph{signed} score.  If its average is
at least $\delta_1>0$, the score is bounded above by a warm constant
$B$ and the set of translations giving score at least
$\delta_1/2$ has measure at least $\delta_1/(2B)$.
Fubini supplies a single translation working on a warm-inverse
fraction of paths.  Translating in the kernel torus preserves the
marked projection.  Repeat at the bounded number of induction
steps.  This produces the asserted single $g^U$ and retains the
original two fields, with only the allotted approximation errors
and target discounts.  Subslices and their full-mass fractions
have warm cost throughout.
\end{proof}

\subsection{Scalar return with a small retained modulus}

\begin{corollary}[Passive scalar return]
\label{ap:passive:scalar-return}
After Proposition~\ref{ap:passive:replacement}, one can select a
modular-good path and its retained warm-cost subslice for which
the forecasting measure $\mathcal K$ of
Lemma~\ref{ap:scalar:measure} has a signed surplus of warm size.
After one further small proportional target discount, its parallel
residue matching can be imposed modulo an integer $M_*$ with warm
logarithmic cost.  The site payoffs still evaluate the original
cutoffs on $g^U(u)$, with the exact old determining conditions.
\end{corollary}

\begin{proof}
Once $g^U$ is fixed, evaluation of the two warm field menus on it
gives menus of bounded site functions with the same supremum
accuracy.  These site functions need not have any regularity.
The scalar comparison, Proposition~\ref{ap:scalar:comparison},
therefore transfers their full-mass signed score to $\mathcal K$
on one modular-good path, including its retained subslice mask.
There are no active flags in this passive step.  The fixed cutoff
functions themselves are recovered after their temporary mesh
approximations, within the prescribed comparison error.

Write $\mathfrak M=\operatorname{lcm}(1,\ldots,Z)$ for the
forecasting modulus.  In this passive use choose $Z$ at warm-pre
cost, after the warm model and slice requirements; auxiliary cold
major comparisons have their own individual-modulus requests and
do not enlarge this $Z$.  The slice strides divide $\mathfrak M$.
Condition on the actual inactive parameters and on the stride
classes.  At a prime $\ell$ which is neither exceptional for
modular decay nor a divisor of the stride, Lemma~\ref{cube:modular-decay}
gives the uniform estimate
\[
                         \|d_\ell-1\|_\infty\leq2\ell^{-2}.
\]
Choose a sufficiently large fixed cutoff $P_0$ so that the product
of these densities over $\ell>P_0$ lies in $1\pm\eta$, with $\eta$
smaller than the allotted proportional discount.  Chinese
remaindering separates these primes from the remaining ones;
mixtures over stride classes cause no change in the uniform
comparison.

Here the comparison is multiplicative separately on the two
nonnegative masses.  If $A_\pm$ denote the masses with all matching
factors and $B_\pm$ those after removing the indicated primes, then
\[
 (1-\eta)B_\pm\leq A_\pm\leq(1+\eta)B_\pm.
\]
Consequently
\[
 B_--\lambda\frac{1-\eta}{1+\eta}B_+
 \geq\frac{A_--\lambda A_+}{1+\eta}.
\]
A fixed proportional discount thus retains a comparable positive
surplus.  In particular $P_0$ need not grow with the inverse of
that possibly tiny surplus.

Let $M_*$ retain the full tested prime powers at primes dividing
the exceptional product $D_0'$, at primes dividing a joint slice
stride $q$, and at primes at most $P_0$.  Every such prime power
has logarithm at most $\log Z$.  Hence
\[
 \log M_*
 \leq\bigl(C_{P_0}+\log_2 D_0'+\log_2 q\bigr)\log Z,
\]
which is warm.  This bounds the full retained depths, not merely
their squarefree product.  The discarded-prime comparison is
pointwise in the site and inactive data, so the exact inactive
equalities and all cutoff conditions are unchanged.  These are
the scalar data needed for the parent-box extraction.
\end{proof}
